\subsection{THE GROUP X(H)}

Let H be a compact Lie group. Denote by X(H) the (commutative) group of continuous homomorphisms from H to the topological group \(C^{*}\) . By Chap. III, §8, no. 1, Th. 1, the elements of X(H) are morphisms of Lie groups; for all \(a \in \mathrm{X}(\mathrm{H})\) , the differential of a is an R-linear map \(\mathrm{L}(a) : \mathrm{L}(\mathrm{H}) \to \mathrm{L}(\mathbf{C}^{*})\) . From now on we identify the Lie algebra of \(C^{*}\) with C in such a way that the exponential map of \(C^{*}\) coincides with the map \(z \mapsto e^{z}\) from C to \(C^{*}\) . Then, to any element \(a \in \mathrm{X}(\mathrm{H})\) is associated an element \(\mathrm{L}(a) \in \mathrm{Hom}_{\mathbf{R}}(\mathrm{L}(\mathrm{H}), \mathbf{C})\) ; we denote by \(\delta(a)\) the element of \(\mathrm{Hom}_{\mathbf{C}}(\mathrm{L}(\mathrm{H})_{(\mathbf{C})}, \mathbf{C})\) that corresponds to it (that is, whose restriction to \(\mathrm{L}(\mathrm{H}) \subset \mathrm{L}(\mathrm{H})_{(\mathbf{C})}\) is equal to \(\mathrm{L}(a)\) ).

For all \(x \in \mathrm{L}(\mathrm{H})\) and all \(a \in \mathrm{X}(\mathrm{H})\) , we have

\[
a (\exp_ {\mathrm{H}} x) = e ^ {\delta (a) (x)},
\]

by functoriality of the exponential map (Chap. III, §6, no. 4, Prop. 10).

We shall often denote the group \(\mathrm{X}(\mathrm{H})\) additively; in that case, we denote the element \(a(g)\) of \(C^{*}\) by \(g^{a}\) . With this notation, we have the formulas

\[
g ^ {a + b} = g ^ {a} g ^ {b}, \quad g \in \mathrm{H}, \quad a, b \in \mathrm{X(H)},
\]

and

\[
(\exp_ {\mathrm{H}} x) ^ {a} = e ^ {\delta (a) (x)}, x \in \mathrm{L} (\mathrm{H}), a \in \mathrm{X} (\mathrm{H}).
\]

Since H is compact, the elements of X(H) take values in the subgroup \(\mathbf{U} = \mathbf{U}(1, \mathbf{C})\) of complex numbers of absolute value 1, so that X(H) can be identified with the group of continuous (or analytic) homomorphisms from H to U. It follows that, for all \(a \in \mathrm{L}(\mathrm{H})\) , the map \(\mathrm{L}(a)\) takes values in the subspace Ri of C, so \(\delta(a)\) maps \(\mathrm{L}(\mathrm{H})\) to Ri.

If H is commutative, \(X(H)\) is simply the (discrete) dual group of H (Spectral Theories, Chap. II, §1, no. 1). If H is commutative and finite, \(X(H)\) can be identified with the dual finite group \(D(H) = \text{Hom}_{\mathbf{Z}}(H, \mathbf{Q}/\mathbf{Z})\) (where, as in Algebra, Chap. VII, §4, no. 9, Example 1, we identify Q/Z with a subgroup of \(C^{*}\) by the homomorphism \(r \mapsto \exp(2\pi ir)\) ).

For any morphism \(f: \mathrm{H} \to \mathrm{H}'\) of compact Lie groups, we denote by \(\mathrm{X}(f)\) the homomorphism \(a \mapsto a \circ f\) from \(\mathrm{X}(\mathrm{H}')\) to \(\mathrm{X}(\mathrm{H})\) . If \(\mathrm{K}\) is a closed normal subgroup of the compact Lie group H, we have an exact sequence of Z-modules \(0 \rightarrow X(H/K) \rightarrow X(H) \rightarrow X(K)\) .

PROPOSITION 1. For any compact Lie group H, the Z-module X(H) is of finite type. It is free if H is connected.

Assume first that H is connected; every element of X(H) vanishes on the derived group D(H) of H, hence we have a homomorphism \(\mathrm{X}(\mathrm{H}/\mathrm{D}(\mathrm{H}))\to\mathrm{X}(\mathrm{H})\) . But \(\mathrm{H}/\mathrm{D}(\mathrm{H})\) is connected and commutative, hence is a torus, and \(\mathrm{X}(\mathrm{H}/\mathrm{D}(\mathrm{H}))\) is a free Z-module of finite type (Spectral Theories, Chap. II, §2, no. 1, Cor. 2 of Prop. 1). In the general case, it follows from the exactness of the sequence

\[
0 \to \mathrm{X} (\mathrm{H} / \mathrm{H} _ {0}) \to \mathrm{X} (\mathrm{H}) \to \mathrm{X} (\mathrm{H} _ {0}),
\]

where \(X(H_{0})\) is free of finite type and \(X(H/H_{0})\) is finite, that \(X(H)\) is of finite type.

PROPOSITION 2. Let H be a commutative compact Lie group, and \((a_{i})_{i\in I}\) a family of elements of X(H); the \(a_{i}\) generate X(H) if and only if the intersection of the Ker \(a_{i}\) reduces to the identity element.

By Spectral Theories, Chap. II, §1, no. 7, Th. 4, the orthogonal complement of the kernel of \(a_i\) is the subgroup \(A_i\) of \(X(H)\) generated by \(a_i\) ; by loc. cit., Cor. 2 of Th. 4, the orthogonal complement of \(\bigcap \operatorname{Ker} a_i\) is the subgroup of \(X(H)\) generated by the \(A_i\) , hence the proposition.

